% Generated draft: requires mathematical and source review.
\documentclass[11pt]{article}
\usepackage[margin=25mm]{geometry}
\usepackage{fontspec}
\setmainfont{Latin Modern Roman}
\usepackage{amsmath,amssymb,mathtools,unicode-math}
\setmathfont{Latin Modern Math}
\usepackage{enumitem,needspace,xurl,hyperref}
\setlength{\parindent}{0pt}
\setlength{\parskip}{5pt}
\setlength{\emergencystretch}{3em}
\newcommand{\N}{\mathbb{N}}
\newcommand{\Z}{\mathbb{Z}}
\newcommand{\Q}{\mathbb{Q}}
\newcommand{\R}{\mathbb{R}}
\newcommand{\C}{\mathbb{C}}
\newcommand{\F}{\mathbb{F}}
\newcommand{\A}{\mathbb{A}}
\newcommand{\PP}{\mathbb{P}}
\newcommand{\E}{\mathbb{E}}
\newcommand{\eps}{\varepsilon}
\newcommand{\dd}{\,\mathrm{d}}
\begin{document}
\section*{UnsolvedMath ID: 11000169}
\textbf{Source problem code:} catalogue entry
\subsection{Definitions and mathematical statement}
Here \(M_g\) denotes the source's genus-\(g\) mapping class group, not a moduli space. The missing preceding construction associates to \(\alpha\in M_g\) a Heegaard gluing class \(\beta\in M_{2g+1}\) for a surface bundle. In the usual closed-fiber setting, the bundle with monodromy \(\alpha\) is the mapping torus
\[
 X_\alpha=\Sigma_g\times[0,1]/(x,1)\sim(\alpha(x),0),
\]
and the source refers to its genus-\(2g+1\) bundle-related splitting. The exact gluing markings and construction are absent.

The compound task is:
\begin{enumerate}
\item Given \(\alpha\), for example as a word in Dehn twists, express the associated \(\beta\) as a corresponding explicit Dehn-twist product.
\item Make precise the suggested implication that equivalent \(\alpha,\alpha'\) yield equivalent associated splittings, and determine whether the converse holds.
\item Recognize whether a given Heegaard splitting is the bundle-related splitting of a fibered \(3\)-manifold.
\item Give the necessary restrictions on a gluing class \(\beta\) that permit reversing the construction and producing the surface bundle.
\end{enumerate}
Equivalence of mapping classes may mean equality of represented elements or conjugacy, while equivalence of splittings may preserve markings, sides, orientation or only the underlying manifold. Those choices are not fixed by the excerpt. A gluing map also depends on handlebody parametrizations, so restrictions on a particular word \(\beta\) must be formulated compatibly with those choices. Recognition of this particular splitting is stronger than merely recognizing that its manifold fibers. No stabilization permission, converse theorem or Dehn-twist formula is supplied.
\subsection{Short English statement}
Make the monodromy-to-Heegaard-gluing construction explicit, determine what equivalence it preserves or detects, and characterize the splittings and gluing maps from which the surface bundle can be recovered.
\subsection{Sources}
\begin{itemize}
\item \emph{Problems on Mapping Class Groups and Related Topics}, edited by Farb, Problem 2.12, supplied PDF page 159. \url{https://www.math.uchicago.edu/~farb/papers/mcgbook.pdf}. Recorded access 2026-07-29; extraction pdf-font-structure. The preceding bundle-related construction and markings are not included; no separate citation is supplied.
\item Local UnsolvedMath registry, \texttt{data/unsolvedmath/problems.json}, permanent ID \(11000169\), exact source code \texttt{catalogue entry}; supplied OPEN-TRIAGE summary reports undetermined status.
\end{itemize}
\end{document}
